% 正文片段：由支持中文的主文档引入；主文档加载 amsmath, amssymb。
% 对应 adaptive_fixed_256 的当前配置；以下是实验设计，不声称已恢复原始生产轨迹。
\subsection{自然节点自适应与固定步长的受迫无滑移流对照}

本实验研究比较相同三阶增量 SDIRK3-mr-ccSAV 格式在联合速度/标量
误差控制和固定步长下的精度、辅助变量及积分成本。
区域为 $\Omega=(0,2\pi)^2$，四周施加齐次无滑移边界条件，
运动黏性系数 $\nu=0.02$，稳态外力为 $f=(-\sin(4y),0)$。
空间上使用 $256\times256$ 的均匀 MAC 网格，终止时间 $T=30$，
取 SAV 阻尼参数 $\gamma=1000$ 和初值 $r_0=0$。
同一物理问题使用张量 Stokes 后端，保留完整动量和散度检查，
残差阈值为 $10^{-9}$；该研究不混用独立固定网格收敛实验的其他外力或初值。

初始速度由壁面夹持流函数的顶点采样经离散旋度生成。
具体地，设置
\begin{equation}
 \psi_0(x,y)=4\sin^2(x/2)\sin^2(y/2)
 \sum_{\substack{k_1,k_2\in\mathbb Z\\0<k_1^2+k_2^2\le100}}
 \frac{[\cos(k_1x)+\sin(k_1x)][\cos(k_2y)+\sin(k_2y)]}
 {(k_1^2+k_2^2)^{3/2}}.
\end{equation}
全部有符号整数模式按确定次序求和，不加随机相位或额外速度归一化。
令 $u_0=\delta_y\psi_0$、$v_0=-\delta_x\psi_0$，
得到离散无散的初始 MAC 速度，壁面仍采用原有无滑移闭合。

自适应成员采用 I 控制器，初始步长为 $10^{-5}$，最小保护步长为 $10^{-12}$，
不设置有限的最大步长上限。安全因子为 $0.9$，常规步长倍率限制在 $[0.5,1.2]$。
速度嵌入差使用绝对容差 $10^{-8}$、相对容差 $2\times10^{-6}$，
标量局部缺陷使用绝对容差 $10^{-4}$。
对于一次候选试步 $h$，指标定义为
\begin{equation}
 E_u=\frac{\|u_{\mathrm{high}}-u_{\mathrm{embedded}}\|_h}
 {10^{-8}+2\times10^{-6}\max(\|u_{\mathrm{old}}\|_h,\|u_{\mathrm{high}}\|_h)},
 \qquad
 E_r=\frac{|r_{\mathrm{trial}}-\mathrm{e}^{-\gamma h}r_{\mathrm{old}}|}{10^{-4}}.
\end{equation}
取 $\max(E_u,E_r)\le1$ 作为接受条件，并保留已有阶段可接受性检查。
速度与标量分别按三阶和二阶局部指标产生步长建议，再采用较小者。
拒绝试步不推进状态或辅助变量；标量容差是局部控制参数，
不是全局 $|r|$ 上界，也不是全局速度误差保证。

当前对照配置为相同三阶 SAV 格式指定八个固定步长
$0.1,0.05,0.025,0.01,0.005,0.0025,0.001,0.0005$。
各组计划进行三次单线程交错重复，积分 CPU 时间排除模型初始化、
Numba/阶段预热、保存及共同参考计算，保留积分内部的求解、检查和诊断成本。
本实验的物理条件和计时边界在组间保持不变，不以不同格式或不同初始化策略替代固定步对照。

按 $0.1$ 间隔提出输出请求，自适应积分不因中间请求截断自然步长；
保存请求附近的最近接受节点，并保留真实物理时刻，误差不在未保存时刻插值。
两个普通 SDIRK3 参考的最大步长分别为 $10^{-4}$ 与 $5\times10^{-5}$，
共同访问实际保存采样时刻的并集。以较细参考计算相对 MAC 速度 $L^2$ 误差，
并要求粗细参考差不超过该成员最大采样误差和终点误差各自的 $5\%$，
才将其参考标记为通过敏感性核验。
该准则是参考分辨率检查而非严格误差界；最大采样误差也不是全部时间节点上的误差上界。
失败成员只保留有效前缀及可用诊断，不将缺失终点填成有效结果。

以上为当前配置的实验设置与比较规则，不是对历史生产轨迹已完成重验的陈述。
当前可访问报告缺少其所引用的原始成员和参考数据，且生产报告仍含参考未完全验证状态。
因此本轮仅提供设置描述，结果 Notebook 等待补回可核验的原始记录；
不以已有汇总统计替代完整时间序列，不据此报告新的速度误差、性能优势或稳定性结论。
